\subsection{Rolle von control.py}
\label{subsec:rolle-control-py}

\path{tools/control.py} ist der öffentliche Einstiegspunkt für die
projektweiten Arbeitsabläufe. Der mit \texttt{argparse} aufgebaute Parser
stellt im untersuchten Stand die Befehle \texttt{init}, \texttt{doctor},
\texttt{install}, \texttt{console}, \texttt{run}, \texttt{stop},
\texttt{test}, \texttt{build}, \texttt{config}, \texttt{db},
\texttt{tauri}, \texttt{container}, \texttt{version}, \texttt{release} und
\texttt{docs} bereit. Gruppenbefehle zeigen ohne Aktion ihre nächste
Befehlsebene; ebenso liefern ein leerer Aufruf und \texttt{--help} ein
Orientierungsmodell, ohne einen Workflow auszuführen
\parencite{kleiveist2026control,kleiveist2026tooling}.

\begin{figure}[H]
  \centering
  \resizebox{\textwidth}{!}{%
  \begin{tikzpicture}[node distance=7mm and 10mm]
    \node[casePrimary,minimum width=40mm,minimum height=18mm] (control) {control.py\\gemeinsamer Einstiegspunkt};

    \node[caseGroup,minimum width=35mm,minimum height=32mm,left=18mm of control] (setupgroup) {};
    \node[font=\sffamily\bfseries\scriptsize,text=caseBlue,above=1mm of setupgroup.north] {Setup};
    \node[caseBox,minimum width=28mm,above=2mm of setupgroup.center] (doctor) {doctor};
    \node[caseBox,minimum width=28mm,below=2mm of setupgroup.center] (init) {init / install};

    \node[caseGroup,minimum width=35mm,minimum height=42mm,right=18mm of control] (devgroup) {};
    \node[font=\sffamily\bfseries\scriptsize,text=caseBlue,above=1mm of devgroup.north] {Entwicklung};
    \node[caseBox,minimum width=28mm,above=7mm of devgroup.center] (run) {run};
    \node[caseBox,minimum width=28mm] at (devgroup.center) (test) {test / build};
    \node[caseBox,minimum width=28mm,below=7mm of devgroup.center] (config) {config};

    \node[caseGroup,minimum width=35mm,minimum height=42mm,below=17mm of control] (specialgroup) {};
    \node[font=\sffamily\bfseries\scriptsize,text=caseBlue,above=1mm of specialgroup.north] {Subsysteme};
    \node[caseBox,minimum width=28mm,above=7mm of specialgroup.center] (db) {db};
    \node[caseBox,minimum width=28mm] at (specialgroup.center) (tauri) {tauri};
    \node[caseBox,minimum width=28mm,below=7mm of specialgroup.center] (release) {release};

    \draw[caseFlow] (control) -- (setupgroup);
    \draw[caseFlow] (control) -- (devgroup);
    \draw[caseFlow] (control) -- (specialgroup);
  \end{tikzpicture}}
  \caption{Befehlsgruppen des zentralen Projekt-Toolings}
  \label{fig:control-command-map}
\end{figure}


Die gemeinsame Oberfläche ersetzt die Fachwerkzeuge nicht. Ihre Handler laden
das aktive Projektmanifest sowie bei Bedarf die Laufzeitkonfiguration und
delegieren anschließend unter anderem an npm und
Vite, Uvicorn und Pytest, Cargo und Tauri, Docker Compose, Alembic oder
PyGitIndex. Damit wird die in Kapitel~\ref{sec:anforderungen} formulierte
Anforderung nach gemeinsamen Entwicklungsbefehlen adressiert, während die
technischen Subsysteme eigenständig bleiben. Die geführte \texttt{console}
ruft dieselben öffentlichen Befehle als Kindprozesse auf und bildet daher eine
alternative Bedienoberfläche, kein zweites Tooling-System.

Erwartbare Bedienfehler werden an der jeweiligen Grenze behandelt: Unbekannte
Parserargumente zeigen Hilfe, Fehlermeldung und nächsten Hilfeschritt und enden
mit Exit-Code~2; fachliche Fehlschläge enden regelmäßig mit Code~1. Hilfe,
erfolgreiche Ausführung und reine Warnzustände liefern Code~0. Bekannte Profil-
und Generierungsfehler werden ohne Traceback ausgegeben. Nur die letzte
Sicherheitsgrenze für unerwartete Exceptions protokolliert zusätzlich den
Stacktrace; ein dort ankommender Abbruch per Tastatur erhält Code~130. Diese
Statuskonvention ermöglicht sowohl interaktive Diagnose als auch die
Auswertung durch CI-Workflows \parencite{kleiveist2026control,kleiveist2026ci}.
